\section{2020-2021学年线性代数II（H）期中答案}

\begin{center}
    任课老师：刘康生\hspace{4em} 考试时长：60分钟
\end{center}
\begin{enumerate}
    \item \begin{enumerate}
        \item 设 $A=\begin{pmatrix}a & b\\c & d\end{pmatrix}$，则
        \[
            T(A)=\begin{pmatrix}a+c & b+d\\a+c & b+d\\0 & 0\end{pmatrix}.
        \]
        因此
        \[
            \im T=\left\{\begin{pmatrix}u & v\\u & v\\0 & 0\end{pmatrix}\mid u,v\in\mathbf{R}\right\},
        \]
        \[
            \ker T=\left\{\begin{pmatrix}a & b\\-a & -b\end{pmatrix}\mid a,b\in\mathbf{R}\right\}.
        \]
        \item 设 $V$ 的自然基与 $W$ 的自然基分别为 $\{v_{11}, v_{12}, v_{21}, v_{22}\}$ 和 $\{w_{11}, w_{12}, w_{21}, w_{22}, w_{31}, w_{32}\}$，则 $T$ 在这两组基下的矩阵为 $M=\begin{pmatrix}1 & 0 & 1 & 0\\0 & 1 & 0 & 1\\1 & 0 & 1 & 0\\0 & 1 & 0 & 1\\0 & 0 & 0 & 0\\0 & 0 & 0 & 0\end{pmatrix}$.

        由于 $r=\dim\im T=2$，设 $\begin{pmatrix}
            E_r  &  O \\ O  &  O
        \end{pmatrix}=P^{-1}MQ^{-1}$，其中 $P$ 和 $Q$ 分别是 $6$ 级和 $4$ 级可逆矩阵，利用行列变换求得矩阵 $P=\begin{pmatrix}
            1 & 0 & 0 & 0 & 0 & 0\\
            0 & 1 & 0 & 0 & 0 & 0\\
            -1 & 0 & 1 & 0 & 0 & 0\\
            0 & -1 & 0 & 1 & 0 & 0\\
            0 & 0 & 0 & 0 & 1 & 0\\
            0 & 0 & 0 & 0 & 0 & 1
        \end{pmatrix}$，$Q=\begin{pmatrix}
            1 & 0 & -1 & 0\\
            0 & 1 & 0 & -1\\
            0 & 0 & 1 & 0\\
            0 & 0 & 0 & 1
        \end{pmatrix}$，$Q^{-1}=\begin{pmatrix}
            1 & 0 & 1 & 0\\
            0 & 1 & 0 & 1\\
            0 & 0 & 1 & 0\\
            0 & 0 & 0 & 1
        \end{pmatrix}$. 由换基公式，所求的 $V$ 的一组基为 $(v_{11}, v_{12}, v_{21}, v_{22})Q^{-1}=(v_{11}, v_{12}, v_{11}+v_{21}, v_{12}+v_{22})$，$W$ 的一组基为 $(w_{11}, w_{12}, w_{21}, w_{22}, w_{31}, w_{32})P=(w_{11}-w_{21}, w_{12}-w_{22}, w_{21}, w_{22}, w_{31}, w_{32})$.
    \end{enumerate}

    \item $f$ 在平面 $U:x_1+x_2+x_3=0$ 上为零，故 $f$ 是 $x_1+x_2+x_3$ 的倍数. 又 $f(1,1,1)=1$，所以
    \[
        f(x_1,x_2,x_3)=\frac{x_1+x_2+x_3}{3}.
    \]

    \item 因 $A^2=A$，对任意 $v\in V$ 有
    \[
        v=Av+(v-Av),
    \]
    其中 $Av\in\im A$，且
    \[
        A(v-Av)=Av-A^2v=0,
    \]
    故 $v-Av\in\ker A$. 又若 $u\in\im A\cap\ker A$，设 $u=Av$，则 $0=Au=A^2v=Av=u$. 因此
    \[
        V=\im A\oplus\ker A.
    \]
    取 $\im A$ 的一组基并接上 $\ker A$ 的一组基，即得可逆矩阵 $P$，使
    \[
        P^{-1}AP=\begin{pmatrix}E_r & O\\O & O\end{pmatrix}.
    \]
\end{enumerate}

\clearpage
